\documentclass[10pt]{article}

\usepackage[utf8]{inputenc}

\usepackage[T1]{fontenc}

\usepackage{amsmath}

\usepackage{amsfonts}

\usepackage{amssymb}

\usepackage[version=4]{mhchem}

\usepackage{stmaryrd}

\usepackage{bbold}



\begin{document}

а характеристич. показатели в остальных особых точках не меняются. С помоцью таких преобразований уравнение (3) приводится к виду



\[

\begin{gathered}

w^{\prime \prime}+\sum_{m=1}^{k} \frac{1-\left(\rho_{2}^{m}+\rho_{1}^{m}\right)}{z-z_{m}} w^{\prime}+ \\

+\left(\bar{\rho}_{1}^{\infty} \rho_{2}^{\infty} z^{n-2}+d_{0} z^{n-3}+\ldots+d_{n-2}\right) \frac{w}{\prod_{m=1}^{k}\left(z-z_{m}\right)}=0 \\

\bar{\rho}_{j}^{\infty}=\rho_{j}^{\infty}+\sum_{m=1}^{k} \rho_{j}^{m}

\end{gathered}

\]



Ф. у. 2-го порядка, имеющее $N$ особых точек, полностью определяется заданием характеристич. шоказателей в этих точках тогда и только тогда, когда $N<4$ m С помощью дробно-линейного преобразования уравнение приводится к виду: а) $N=1, \tilde{w}^{\prime \prime}=0$; б) $N=2$, $\zeta^{2} \tilde{w}^{\prime \prime}+a \zeta \tilde{w}^{\prime}+b \tilde{w}=0 \quad$ (Э й ле р а у р а в не в и е); в) $N=3$ - Папперица уравнение (или уравнение Римана).



Матричное Ф. у. имеет вид



\[

\frac{d W}{d z}=\sum_{m=1}^{k} \frac{W U_{m}}{z-z_{m}}

\]



где $z_{1}, \ldots, z_{k}-$ различные точки, $W-$ матрица-функция порядка $n \times n, U_{m} \neq 0$-постоянные матрицы. Матрица $U_{m}$ наз. Дифф е р енциа ль ной под с т а но вкой в точке $z_{m}$. Пусть $\gamma$-простая замкнутая кривая с началом в неособой точке $b$, положительно ор иентированная и содержащая внутри себя только одну особую точку $z_{m}$. Если $W(z)$ - голоморфное в тотке $b$ решение уравнения (4), то при аналитич. продолжении вдоль $\gamma$



\[

W \rightarrow V_{m} W

\]



где $V_{m}$ - постоянная матрица, наз. и н т е г р а л ьно й ІІ о д с т а по о к о й в $z_{m}$. А. Пуанкаре (H. Poincaré, см. [2]) поставил для систем вида (4) задачу, к-рая наз. І р я мой регу ля р ной д а ч ө й П у а н к а р е. Она состоит из следующих трех задач:



А) представление решения $W(z)$ во всей области его существования;



Б) построение интегральных подстановок в точках $z_{m} ;$



В) аналитич. характеристика особенностей решений.



В частности, решение задачи Б) позволяет построить групшу монодромии уравнения (4). Решение задачи Пуанкаре было получево И. А. Лаппо-Данилевским [3]. Пусть $L_{b}\left(z_{j_{1}}, \ldots,\left.z_{j_{v}}\right|^{z}\right), j_{m} \in\{1, \ldots, k\}, v=1,2, \ldots$, - гиперлогарифмы:



\[

\begin{aligned}

L_{b}\left(z_{m} \mid z\right) & =\int_{b}^{z} \frac{d z}{z-z_{m}}, L_{b}\left(z_{j_{1}}, \ldots, z_{j_{v}} \mid z\right)= \\

& =\int_{b}^{z} \frac{L_{b}\left(z_{j_{1}}, \ldots,\left.z_{j v-1}\right|^{z}\right)}{z-z_{j v}} d z,

\end{aligned}

\]



$W_{0}(z)$ - элемент (росток) в точке $b$ решения уравнения (4), нормированный условием $W_{0}(b)=I$ и $W(z)-$ аналитическая в области $D$ матрица-функция, порожденная этим әлементом. Тогда $W(z)$ есть целая функция от матриц $U_{1}, \ldots, U_{k}$ и разлагается в ряд



$W(z)=I+\sum_{v=1}^{\infty} \sum_{j_{1}, \ldots, j_{v}}^{(1, \ldots, k)} U_{j_{1}} \ldots U_{j v} L_{b}\left(z_{j_{1}}, \ldots, z_{j_{v}} z\right)$, к-рый сходится равномерно по $z$ на любом компакте $K \subset D$. Интегральная подстановка $V_{m}$ в точке $z_{m}$, отвечающая решению $W(z)$, есть целая функция от матриц $U_{1}, \ldots, U_{k}$ и разлагается в ряд



$V_{m}=I+\sum_{v=1}^{\infty} \sum_{j_{1}, \ldots, j_{v}}^{(1, \ldots, k)} U_{j_{1}} \ldots U_{j v} P_{j}\left(z_{j_{1}}, \ldots, z_{j v} \mid b\right)$, где $P_{j}$ выражаются через гиперлогарифмы (см. [3], $[6])$



Получены такжке формулы, дающие решение задачи B) (см. [3]).



Jum.: [1] F u c h s L., «J. reine und angew. Math.», 1866, Bd 66, S. $121-60 ; 1868$, Bd 68, S. 354-85, [2] II у а н к а р е А., Избр. труды, пер. с франц., т. 3, М., 1974; [3] Јі а п п о-Д, а н и л е вс к и й И. А., Применение функций от матриц к теории линей-

ных систем обыкновенных дифференциальных уравнений, пер. с франц., М.,1957; [4] К о д д и н г т о н Э. А.,, Л е в и н с о н Н., Теория обыкновенных дифференциальных уравнений, пер. с англ., М., 1958, [5] Г о л у б е в В. В., Лекции по аналитической теории дифференциальных уравнений, 2 изд., М.-Ј., 1950; [6] С м и р н о в В. И., Курс высшей математики, т. 3, ч. 2, уравнения, пер. с англ., Хар., 1939.



ФУКСОВА ГРУППА - әискретная группа голоморфных преобразований (открытого) круга $K$ на сфере Римана, т. е. круга или полуплоскости на комплексной плоскости. Чаще всего в качестве $K$ берут верхнюю полуплоскость



\[

U=\{z \in \mathbb{C} \mid \operatorname{Im} z>0\}

\]



или единичный круг



\[

D=\{z \in \mathbb{C}|| z \mid<1\}

\]



В первом случае элементы Ф. Г. являются дробнолинейными преобразованиями



\[

z \longmapsto \frac{a z+b}{c z+d}

\]



с действительными коәффициентами, и Ф. г. представляет собой не что иное, как дискретную подгрупnу группы $P S L_{2}$. Во втором случае элементы Ф. г. являются дробно-линейными преобразованиями с псевдоунитарными матрицами.



Если рассматривать круг $K$ как конформную модель плоскости Лобачевского, то Ф. г. может быть определена как дискретная группа его движений, сохраняющих ориентацию. Ф. Г. представляют собой частный случай клейновых әрупn.



Произвольные $\Phi$. г. впервые рассматривались А. Пуанкаре (Н. Poincaré, см. [2]) в 1882 в связи с проблемой униформизации. Группы были названы им фукковыми в честь Л. Фукса, работа [1] к-рого стимулировала введение этого понятия. ДЛя описания Ф. Г. Пуанкаре применил комбинаторно-геометрич. метод, ставший впоследствии одним из основных методов теории дискретных грушп преобразований. Понятие Ф. г. послужило основой для теории автоморфных функций, созданной А. Пуанкаре и Ф. Клейном (F. Klein).



Ф. г., сохраняющая какую-либо точку в замыкании $\bar{K}$ круга $K$ или прямую в смысле геометрии Лобачевского, наз. э л е м е в т а р н о й. Если $Г-$ неэлементарная Ф. г., то множество $L(\Gamma)$ предельных точек орбиты точки $x \in \bar{K}$, лежащее на граничной окружности $\partial K$, не зависит от $x$ и наз. І $\mathrm{p}$ е д е л ь н ы м м н ож е с т в о м г р у п п ы $\Gamma$. Группа $\Gamma$ наз. Ф. г. 1-г о р о д а, если $L(\Gamma)=\partial K$, и $2-\Gamma \circ$ р о д а- в противном случае (тогда $L(\Gamma)$ - нигда не плотное совершенное подмножество в $\partial K)$.



Конечно порожденная Ф. г. является Ф. Г. 1-го рода тогда и только тогда, когда плошадь (в смысле геометрии Лобачевского) ее фундаментальной области конечна. В качестве фундаментальной области такой группы Г всегда может быть выбран выпуклый многоугольник $\boldsymbol{P}$ плоскости Лобачевского со сторонами



$a_{1}, b_{1}^{\prime}, a_{1}^{\prime}, b_{1}, \ldots, a_{g}, b_{g}^{\prime}, a_{g}^{\prime}, b_{g}, c_{1}, c_{1}^{\prime}, \ldots, c_{n}, c_{n}^{\prime}$ таким образом, что



\[

a_{i}=\alpha_{i}\left(a_{i}^{\prime}\right), \quad b_{i}=\beta_{i}\left(b_{i}^{\prime}\right), \quad c_{i}=\gamma_{i}\left(c_{i}^{\prime}\right)

\]



для нек-рых элементов



\[

\alpha_{1}, \ldots, \alpha_{g}, \beta_{1}, \ldots, \beta_{g}, \gamma_{1}, \ldots, \gamma_{n},

\]





\end{document}